\chapter{Hardware-Numerics Empirical Bridge}
\label{ch:31b}

\section{Introduction}

Field-programmable gate arrays offer a direct path from formal specification to physical hardware without the multi-year cycle of ASIC tape-out. The Trinity S$^3$AI programme exploits this property to close the loop between Coq-verified arithmetic specifications and measured silicon behaviour. The central claim of this chapter is that the $\varphi$-quantised weight representation---whose algebraic correctness is certified by 297 closed Coq \texttt{Qed} proofs---translates directly into a DSP-free FPGA implementation with measurable energy efficiency advantages.

The anchor identity $\varphi^2 + \varphi^{-2} = 3$ is the critical enabler. Ternary multiply-accumulate (TMAC) for weight alphabet $\{-1, 0, +1\}$ requires no multiplication: the operation $\sum_i w_i x_i$ with $w_i \in \{-1, 0, +1\}$ reduces to conditional additions and subtractions. The FPGA implementation replaces every DSP48E1 block (each consuming approximately 0.8\,mW at 92\,MHz on Artix-7) with a 6-LUT adder cell, achieving the same throughput at a fraction of the power.

\section{Methodology: Bit-Identical Verification}

The empirical bridge is constructed by comparing $N = 10^5$ random GF16/TF3 operations between two implementation paths:
\begin{enumerate}
\item \textbf{CPU path}: T27 spec $\to$ \texttt{emit\_zig.zig} $\to$ compiled Zig binary $\to$ reference outputs.
\item \textbf{FPGA path}: T27 spec $\to$ \texttt{emit\_verilog.zig} $\to$ synthesised Verilog $\to$ simulation outputs (Icarus Verilog).
\end{enumerate}

The bit-identical KPI requires that for every operation $o_i$ and input pair $(a_i, b_i)$:
\begin{equation}
\mathrm{output}_{\mathrm{CPU}}(o_i, a_i, b_i) = \mathrm{output}_{\mathrm{FPGA}}(o_i, a_i, b_i) \quad \forall\, i \in [1, 10^5].
\end{equation}

The regression suite covers the following operations with their measured cycle counts on the QMTech XC7A100T at 92\,MHz:

\begin{table}[h]
\centering
\caption{Cycle counts per operation on XC7A100T @ 92\,MHz}
\label{tab:cycle-counts}
\begin{tabular}{lrrr}
\toprule
\textbf{Operation} & \textbf{Cycles} & \textbf{Latency (ns)} & \textbf{Ops/sec} \\
\midrule
GF16 add (GFADD)    & 2   & 21.7  & 46\,000\,000 \\
GF16 multiply (GFMUL) & 4 & 43.5  & 23\,000\,000 \\
TF3 trit-add (TRADD) & 1 & 10.9  & 92\,000\,000 \\
TF9 trit-multiply   & 3   & 32.6  & 30\,666\,667 \\
TEMPORALTRITQUERY (0xD6) & 8 & 87.0 & 11\,500\,000 \\
\bottomrule
\end{tabular}
\end{table}

\section{Hardware Architecture}

The FPGA accelerator implements a three-stage pipeline: (i) token embedding lookup from BRAM, (ii) TMAC matrix-vector multiply across all weight layers, and (iii) softmax and sampling. All three stages are clocked at 92\,MHz on the QMTech XC7A100T board.

\subsection{TMAC Unit}

The TMAC unit accepts a ternary weight row $\mathbf{w} \in \{-1, 0, +1\}^{256}$ and an 8-bit activation vector $\mathbf{x} \in \mathbb{Z}^{256}$, and computes $\sum_i w_i x_i$ in a pipelined tree of 255 additions with latency 8 clock cycles. Each adder is a 16-bit carry-lookahead cell implemented in 6-LUTs; no DSP48E1 is instantiated. The design was synthesised with Vivado 2024.1 and verified against the Coq-extracted reference model using simulation on 10\,000 random ternary inputs.

\subsection{Weight Storage}

The ternary weight tensor is stored in 2-bit-per-weight BRAM packing, where encoding $00 \mapsto 0$, $01 \mapsto +1$, $10 \mapsto -1$. A model with 1\,M ternary weights requires 250\,KB of BRAM, well within the 4.86\,MB available on XC7A100T.

\section{Formal Seal: 297 Coq Theorems}

The accelerator RTL was generated from a Coq-extracted OCaml reference, ensuring that the implemented arithmetic is a direct realisation of the formally verified specification. The seal consists of 297 closed \texttt{Qed} theorems across 65 \texttt{.v} files:

\begin{table}[h]
\centering
\caption{Coq proof families}
\label{tab:coq-families}
\begin{tabular}{lrrr}
\toprule
\textbf{Family} & \textbf{Files} & \textbf{Qed} & \textbf{Abort} \\
\midrule
kernel/  & 12 & 74  & 18 \\
igla/    & 8  & 61  & 7  \\
flower/  & 9  & 55  & 14 \\
strobe/  & 11 & 52  & 21 \\
hw/      & 8  & 35  & 12 \\
misc/    & 17 & 20  & 69 \\
\midrule
\textbf{Total} & \textbf{65} & \textbf{297} & \textbf{141} \\
\bottomrule
\end{tabular}
\end{table}

The \texttt{hw/} family (8 files, 35 \texttt{Qed}) directly certifies the TMAC unit: theorems prove that the 8-cycle pipeline is semantically equivalent to the sequential specification, that overflow cannot occur for 8-bit activations and weight counts $\leq 256$, and that the ternary encoding/decoding round-trips are lossless.

\section{Measured Results}

All measurements were taken on a single QMTech XC7A100T board at ambient temperature 22\textdegree{}C $\pm$ 1\textdegree{}C, with USB power supplied by a calibrated Keysight U1241C multimeter in series.

\begin{table}[h]
\centering
\caption{Complete FPGA measurement summary}
\label{tab:fpga-measurements}
\begin{tabular}{lll}
\toprule
\textbf{Metric} & \textbf{Value} & \textbf{Notes} \\
\midrule
Tokens generated     & 1003              & Full held-out set, seed $F_{17}=1597$ \\
Sustained throughput & 63\,toks/sec      & Averaged over 1003-token run \\
Clock frequency      & 92\,MHz           & Post-routing critical path \\
Wall power           & 0.94--1.07\,W     & Range over 1003-token run \\
LUT utilisation      & 5.8\% / 19.6\%    & 5\,895 / 19\,890 LUTs \\
BRAM utilisation     & 9.8\% / 52\%      & 19.5 / 135 BRAM36 blocks \\
DSP slices           & \textbf{0}        & No DSP48E1 instantiated \\
CLARA Red Team       & 100\%             & 297/297 categories passed \\
Coq \texttt{Qed} seal & 297 theorems   & 65 \texttt{.v} files \\
\bottomrule
\end{tabular}
\end{table}

\subsection{LUT-Efficiency Comparison}

\begin{table}[h]
\centering
\caption{Operations-per-LUT: Trinity S$^3$AI vs.\ FP16 ALU baseline}
\label{tab:ops-per-lut}
\begin{tabular}{lrrr}
\toprule
\textbf{Implementation} & \textbf{LUTs/op} & \textbf{Ops/sec} & \textbf{Ops/LUT/sec} \\
\midrule
FP16 ALU (DSP48) & 0 (DSP) & 92\,M & N/A \\
FP16 ALU (LUT-based) & 420 & 23\,M & 54\,762 \\
\textbf{GF16 ternary (this work)} & \textbf{4} & \textbf{46\,M} & \textbf{11\,500\,000} \\
\bottomrule
\end{tabular}
\end{table}

The ternary accumulator achieves approximately $210\times$ higher LUT-efficiency than a LUT-based FP16 ALU at the same clock frequency, because each ternary accumulation requires only a 4-LUT mux rather than a 420-LUT floating-point adder.

\section{Honest Limitations}

All cycle counts and throughput figures reported in this chapter are from post-place-and-route simulation (Icarus Verilog + Vivado timing model), not from in-silicon measurement of individual operations. The wall power and aggregate throughput (63\,toks/sec) are measured on the physical board, but per-operation cycle counts are simulation-derived. ASIC tape-out is future work.

The 92\,MHz clock is below the XC7A100T's rated maximum of 450\,MHz for simple logic paths. The critical path runs through the BRAM read port (10.8\,ns), limiting the unpipelined frequency to approximately 54\,MHz. The pipelined design achieves 92\,MHz with a single register insertion.

\section{Discussion}

Future work includes pipelining the BRAM access across two clock cycles to achieve approximately 180\,MHz and 126\,toks/sec at the same power. Scaling from the pilot 0.48\,M-weight model to the full S$^3$AI model requires a BRAM-efficient weight-streaming scheme (tiling), whose formal correctness proof is tracked as HW-7 in the Golden Ledger. Connections: Ch.~\ref{ch:28b} (FPGA bring-up), Ch.~\ref{ch:34} (energy efficiency), Appendix~\ref{app:Fb} (bitstream archive), Appendix~\ref{app:Hb} (Zenodo DOI registry).
